\documentclass[11pt,a4paper,oneside]{report}

\usepackage{notes}

% LighTex:TitleMetadata

\begin{document}

% LighTex:TitlePage

% LighTex:ChapterContentStart
\chapter{Window Functions}

Window functions let us study a finite portion of a signal while controlling
the artifacts introduced at its boundaries. They appear in spectral analysis,
filter design, statistics, and numerical computation.

\begin{definition}[Window function]
A \term{window function} is a finite sequence
\[
  w[0],w[1],\ldots,w[N-1]
\]
used to weight a signal before it is analysed. The windowed signal is
\[
  x_w[n] = x[n]w[n].
\]
\end{definition}

\begin{theorem}[Symmetry]
If a real-valued window satisfies \(w[n]=w[N-1-n]\), then its discrete-time
Fourier transform has linear phase.
\end{theorem}

\begin{proof}
Pair the terms indexed by \(n\) and \(N-1-n\) in the Fourier sum. Their common
phase factor is \(e^{-i\omega(N-1)/2}\), while the remaining factor is real and
even in \(\omega\).
\end{proof}

\section{The Hann window}

The Hann window is defined by
\[
  w[n] =
  \frac{1}{2}\left(1-\cos\frac{2\pi n}{N-1}\right),
  \qquad 0\le n<N.
\]
Its samples taper smoothly to zero, reducing spectral leakage at the cost of a
wider main lobe.

\begin{keyidea}
Window selection is a trade-off: a narrow main lobe improves frequency
resolution, while low side lobes suppress leakage from strong neighbouring
components.
\end{keyidea}

\begin{example}
For \(N=8\), evaluating the formula gives a symmetric sequence:
\[
  (0,\;0.188,\;0.611,\;0.950,\;0.950,\;0.611,\;0.188,\;0).
\]
\end{example}

% LighTex:CodeExample
% LighTex:ChapterContentEnd

\end{document}
